\section{2016 Algebra \& Number Theory}

\subsection*{Individual}
\begin{exercise}
Problem 1 (20 points). Let $E$ be a linear space over $\mathbb{R}$, of finite dimension $n \geq 2$, equipped with a positive definite symmetric bilinear form $\langle \cdot, \cdot \rangle$. Let $u_1, u_2, \ldots, u_n$ be a basis of $E$. Let $v_1, v_2, \ldots, v_n$ be the dual basis, that is,
$$\langle u_i, v_j \rangle = \begin{cases} 1 & \text{if } i = j, \\ 0 & \text{if } i \neq j, \end{cases}$$
for all $i, j = 1, 2, \ldots, n$.

(a) (8 points) Assume that $\langle u_i, u_j \rangle \leq 0$ for all $1 \leq i < j \leq n$. Show that there is an orthogonal basis $u_1', u_2', \ldots, u_n' of E$ such that $u_i'$ is a non-negative linear combination of $u_1, u_2, \ldots, u_i$, for all $i = 1, 2, \ldots, n$.

(b) (6 points) With the same assumption as in Part (a), show that $\langle v_i, v_j \rangle \geq 0$ for all $1 \leq i < j \leq n$.

(c) (6 points) Assume that $n \geq 3$. Show that the condition $\langle u_i, u_j \rangle \geq 0$ for all $1 \leq i < j \leq n$ does not imply that $\langle v_i, v_j \rangle \leq 0$ for all $1 \leq i < j \leq n$.
\end{exercise}
\begin{solution}
\textbf{Motivation and Strategy:} 
This problem explores the relationship between a basis and its dual basis in an inner product space, particularly how sign conditions on the inner products of the basis vectors translate to the dual basis. 
\begin{itemize}
    \item Part (a) asks us to preserve "positivity" while orthogonalizing. The standard tool for orthogonalization is the Gram-Schmidt process. We must verify that the coefficients produced by the process remain non-negative under the given constraints.
    \item Part (b) connects the dual basis to the inverse of the Gram matrix. The condition on the Gram matrix (non-positive off-diagonal entries) characterizes it as a \textit{Stieltjes matrix}, a subclass of \textit{M-matrices}, which are known to have non-negative inverses.
    \item Part (c) requires a counterexample. We seek a configuration where vectors are "close" (positive inner products) but their duals do not satisfy the opposite sign condition.
\end{itemize}

(a) We use the Gram-Schmidt process to construct the orthogonal basis. Let $u_1' = u_1$. For $i = 2, \dots, n$, define recursively:
$$u_i' = u_i - \sum_{j=1}^{i-1} \lambda_{ij} u_j', \quad \text{where } \lambda_{ij} = \frac{\langle u_i, u_j' \rangle}{\langle u_j', u_j' \rangle}.$$
By construction, $\langle u_i', u_k' \rangle = 0$ for all $k < i$. We prove by induction on $i$ that $u_i'$ is a non-negative linear combination of $u_1, \dots, u_i$, which also implies $\langle u_i, u_j' \rangle \leq 0$ for $j < i$.

\textit{Base Case:} For $i=1$, $u_1' = u_1$ is trivially a non-negative linear combination of $u_1$.

\textit{Inductive Step:} Assume the claim holds for indices $1, \dots, i-1$. For any $j < i$, $u_j' = \sum_{k=1}^j c_{jk} u_k$ with $c_{jk} \geq 0$. Then:
$$\langle u_i, u_j' \rangle = \sum_{k=1}^j c_{jk} \langle u_i, u_k \rangle.$$
Since $k \leq j < i$, the problem statement assumes $\langle u_i, u_k \rangle \leq 0$. Since $c_{jk} \geq 0$, the sum $\langle u_i, u_j' \rangle \leq 0$.
Consequently, the Gram-Schmidt coefficients satisfy $\lambda_{ij} = \frac{\langle u_i, u_j' \rangle}{\|u_j'\|^2} \leq 0$.
Rewriting the recursive step as $u_i' = u_i + \sum_{j=1}^{i-1} (-\lambda_{ij}) u_j'$, we see that $u_i'$ is a non-negative linear combination of $u_i$ and $\{u_j'\}_{j<i}$. Since each $u_j'$ is already a non-negative combination of $\{u_1, \dots, u_j\}$, $u_i'$ is a non-negative combination of $\{u_1, \dots, u_i\}$.

(b) Let $G$ be the Gram matrix of the basis $\{u_i\}$, where $G_{ij} = \langle u_i, u_j \rangle$. By assumption, $G$ is symmetric positive definite and $G_{ij} \leq 0$ for $i \neq j$. Such a matrix is called a \textbf{Stieltjes matrix}.
The dual basis $\{v_j\}$ is defined by $\langle u_i, v_j \rangle = \delta_{ij}$. If we write $v_j = \sum_k w_{jk} u_k$, the condition becomes $\sum_k w_{jk} \langle u_k, u_i \rangle = \delta_{ji}$, which in matrix form is $W G = I$, where $W = (w_{ji})$. Thus, the matrix of inner products of the dual basis is:
$$\langle v_i, v_j \rangle = \langle \sum_k w_{ik} u_k, \sum_l w_{jl} u_l \rangle = \sum_{k,l} w_{ik} G_{kl} w_{jl} = (W G W^T)_{ij} = (G^{-1})_{ij}.$$
A key theorem in matrix theory states that if $G$ is an M-matrix (specifically a Stieltjes matrix here), then $G^{-1}$ has entries that are all non-negative.
Therefore, $(G^{-1})_{ij} = \langle v_i, v_j \rangle \geq 0$ for all $i, j$, proving the claim.

(c) We show by counterexample for $n=3$ that $\langle u_i, u_j \rangle \geq 0$ does not imply $\langle v_i, v_j \rangle \leq 0$.
Consider the Gram matrix $G = \begin{pmatrix} 1 & 0.5 & 0.5 \\ 0.5 & 1 & 0.1 \\ 0.5 & 0.1 & 1 \end{pmatrix}$. 
1. \textbf{Positive Definiteness:} $\det(G) = 0.54 > 0$.
2. \textbf{Inner Products:} All $G_{ij} \geq 0$, so $\langle u_i, u_j \rangle \geq 0$.
3. \textbf{Dual Inner Products:} $(G^{-1})_{23} = \frac{-(-0.15)}{0.54} = \frac{0.15}{0.54} > 0$.
Since $\langle v_2, v_3 \rangle$ is positive, the implication fails.
\end{solution}

\begin{exercise}
Problem 2 (20 points). Let $d \geq 1$ and $n \geq 1$ be integers.

(a) (5 points) Show that there are only finitely many subgroups $G \subseteq \mathbb{Z}^d$ of index $n$. Let $f_d(n)$ denote the number of such subgroups.

(b) (5 points) Let $g_d(n)$ denote the number of subgroups $H \subseteq \mathbb{Z}^d$ of index $n$ such that the quotient group is cyclic. Show that $f_d(mn) = f_d(m)f_d(n)$ and $g_d(mn) = g_d(m)g_d(n)$ for coprime positive integers $m$ and $n$.

(c) (5 points) Compute $g_d(p^r)$ for every prime power $p^r$, $r \geq 1$.

(d) (5 points) Compute $f_2(20)$.
\end{exercise}
\begin{solution}
\textbf{Motivation and Strategy:} 
This problem deals with the classification and counting of subgroups of the free abelian group $\mathbb{Z}^d$. 
\begin{itemize}
    \item Part (a) uses the \textbf{Hermite Normal Form (HNF)}, which provides a unique basis for any sublattice of $\mathbb{Z}^d$. The "size" of the subgroup is its index, which corresponds to the volume of the fundamental parallelotope (the determinant of the basis matrix).
    \item Part (b) explores multiplicativity. Since index is multiplicative for coprime factors, we can leverage the structure of the quotient groups $\mathbb{Z}^d/K$ using the Chinese Remainder Theorem.
    \item Part (c) counts cyclic quotients. This is equivalent to counting surjective homomorphisms to $\mathbb{Z}/p^r\mathbb{Z}$ up to automorphisms of the target.
    \item Part (d) applies the theory to a specific case. For $d=2$, the HNF matrices are easy to enumerate.
\end{itemize}

(a) Every subgroup $G \subseteq \mathbb{Z}^d$ of index $n$ is a sublattice of rank $d$. Such a subgroup can be uniquely represented by a matrix $M$ in \textbf{Hermite Normal Form} whose rows form a basis for $G$:
$$M = \begin{pmatrix} m_{11} & 0 & \cdots & 0 \\ m_{21} & m_{22} & \cdots & 0 \\ \vdots & \vdots & \ddots & \vdots \\ m_{d1} & m_{d2} & \cdots & m_{dd} \end{pmatrix}$$
where $m_{ii} > 0$, $0 \leq m_{ij} < m_{ii}$ for $j < i$, and the index is given by $[\mathbb{Z}^d : G] = \det(M) = \prod_{i=1}^d m_{ii} = n$.
Since there are finitely many ways to factor $n$ into $d$ positive integers, and for each factor $m_{ii}$ there are finitely many choices for off-diagonal entries, there are only finitely many such matrices, and thus $f_d(n)$ is finite.

(b) \textbf{Multiplicativity:} Let $n = mk$ with $\gcd(m, k) = 1$. A subgroup $G \subseteq \mathbb{Z}^d$ of index $mk$ corresponds to a quotient $Q = \mathbb{Z}^d/G$ of order $mk$. By the Primary Decomposition Theorem for finite abelian groups, $Q \cong Q_m \oplus Q_k$ where $|Q_m|=m$ and $|Q_k|=k$. This established a bijection $Sub_{mk}(\mathbb{Z}^d) \cong Sub_m(\mathbb{Z}^d) \times Sub_k(\mathbb{Z}^d)$, so $f_d(mk) = f_d(m)f_d(k)$. The same argument holds for $g_d(n)$.

(c) A subgroup $H \subseteq \mathbb{Z}^d$ has $\mathbb{Z}^d/H \cong \mathbb{Z}/p^r\mathbb{Z}$ if and only if $H$ is the kernel of a surjective homomorphism $\phi: \mathbb{Z}^d \to \mathbb{Z}/p^r\mathbb{Z}$. 
The number of surjective homomorphisms is $p^{rd} - p^{(r-1)d} = p^{(r-1)d}(p^d - 1)$.
Two surjections have the same kernel if and only if they differ by an automorphism of $\mathbb{Z}/p^r\mathbb{Z}$. The number of such automorphisms is $|\mathrm{Aut}(\mathbb{Z}/p^r\mathbb{Z})| = \phi(p^r) = p^{r-1}(p-1)$.
Therefore:
$$g_d(p^r) = \frac{p^{(r-1)d}(p^d - 1)}{p^{r-1}(p-1)} = p^{(r-1)(d-1)} \frac{p^d - 1}{p-1} = p^{(r-1)(d-1)} \sum_{j=0}^{d-1} p^j.$$

(d) For $d=2$, $f_2(n) = \sum_{a|n} a = \sigma_1(n)$.
For $n = 20 = 2^2 \cdot 5^1$:
$$f_2(20) = \sigma_1(4)\sigma_1(5) = (1+2+4)(1+5) = 7 \cdot 6 = 42.$$
\end{solution}

\begin{exercise}
Problem 3 (20 points). Let $A$ be a complex $m \times m$ matrix. Assume that there exists an integer $N \geq 0$ such that $t_n = \mathrm{tr}(A^n)$ is an algebraic integer for all $n \geq N$. The goal of this problem is to show that the eigenvalues $a_1, \ldots, a_m$ of $A$ are algebraic integers.

(a) (10 points) Show that there exist algebraic numbers $b_{ij} \in \mathbb{C}$, $1 \leq i, j \leq m$ such that
$$a_i^n = \sum_{j=1}^m b_{ij} t_{n+j-1}$$
for all $n \geq 0$ and all $1 \leq i \leq m$. In particular, $a_1, \ldots, a_m$ are algebraic numbers.

(b) (8 points) Let $R$ be the ring of all algebraic integers in $\mathbb{C}$ and let $K$ be the field of all algebraic numbers in $\mathbb{C}$. Show that for $a \in K$, if $R[a]$ is contained in a finitely-generated $R$-submodule of $K$, then $a \in R$.

(c) (2 points) Conclude that $a_1, \ldots, a_m$ are algebraic integers.
\end{exercise}
\begin{solution}
\textbf{Motivation and Strategy:} 
This problem links the traces of matrix powers (Newton sums) to the eigenvalues of the matrix. 
\begin{itemize}
    \item Part (a) uses the relation between Newton sums and eigenvalues.
    \item Part (b) focuses on the property of \textbf{integral extensions}.
    \item Part (c) combines these to show that if traces are eventually algebraic integers, the eigenvalues must be as well.
\end{itemize}

(a) Let $a_1, \ldots, a_m$ be the eigenvalues. $t_n = \sum_{i=1}^m a_i^n$.
The characteristic polynomial $P(x) = \prod (x - a_i)$ has coefficients in $K$, so $a_i \in K$.
One can show $a_i^n = \sum_j b_{ij} t_{n+j-1}$ for some $b_{ij} \in K$ by using Jordan form or Vandermonde-type relations.

(b) If $R[a] \subseteq M$, where $M$ is a finitely generated $R$-module, then $a$ is integral over $R$. Since $R$ is integrally closed, $a \in R$.

(c) For $n \geq N$, $t_{n+j-1} \in R$. Thus $a_i^n$ is integral over $R$. This implies $a_i$ is integral over $R$, hence $a_i \in R$.
\end{solution}

\begin{exercise}
Problem 4 (20 points). Let $E$ be a Euclidean plane. For each line $l$ in $E$, write $s_l \in \mathrm{Iso}(E)$ for the reflection with respect to $l$, where $\mathrm{Iso}(E)$ denotes the group of distance-preserving bijections from $E$ to itself.

(a) (6 points) Let $l_1$ and $l_2$ be two distinct lines in $E$. Find the necessary and sufficient condition that $s_{l_1}$ and $s_{l_2}$ generate a finite group.

(b) (7 points) Let $l_1, l_2$ and $l_3$ be three pairwise distinct lines in $E$. Assume that $s_{l_1}, s_{l_2}$ and $s_{l_3}$ generate a finite group. Show that $l_1, l_2, l_3$ intersect at a point.

(c) (7 points) Let $G$ be a finite subgroup of $\mathrm{Iso}(E)$ generated by reflections. Show that $G$ is generated by at most two reflections.
\end{exercise}
\begin{solution}
(a) $s_{l_1} s_{l_2}$ is a translation (if $l_1 \parallel l_2$) or a rotation by $2\theta$. A translation generates an infinite group. A rotation generates a finite group iff $2\theta \in 2\pi \mathbb{Q}$, i.e., $\theta \in \pi \mathbb{Q}$.

(b) Every finite group of isometries of the Euclidean plane has a fixed point $P$. Then $s_i(P) = P$ for each $i=1, 2, 3$, which means $P \in l_i$.

(c) Since all reflections in $G$ pass through $P$, $G \subset O(2)$. The only finite subgroups of $O(2)$ generated by reflections are dihedral groups $D_{2n}$, which are generated by two reflections.
\end{solution}

\begin{exercise}
Problem 5 (20 points). Let $G$ be a finite group of order $2^n m$ where $n \geq 1$ and $m$ is an odd integer. Assume that $G$ has an element of order $2^n$. The goal of this problem is to show that $G$ has a normal subgroup of order $m$.

(a) (5 points) Show that if $M$ is a normal subgroup of $G$ of order $m$, then $M$ is the only subgroup of $G$ of order $m$.

(b) (5 points) Let $N$ be a normal subgroup of $G$ and let $P$ be a 2-Sylow subgroup of $G$. Show that $P \cap N$ is a 2-Sylow subgroup of $N$.

(c) (5 points) Show that the homomorphism $G \to \{\pm 1\}$ carrying $g$ to $\mathrm{sgn}(l_g)$ is surjective. Here $\mathrm{sgn}(l_g)$ denotes the sign of the permutation $l_g: G \to G$ given by left multiplication by $g$.

(d) (5 points) Deduce by induction on $n$ that $G$ has a normal subgroup of order $m$.
\end{exercise}
\begin{solution}
(a) Let $M$ be a normal subgroup of order $m$ and $M'$ any subgroup of order $m$. $M M'$ is a subgroup of odd order, so $|M M'| \mid m$, implying $M' \subseteq M$.

(b) $[N : P \cap N] = [P N : P]$, which is odd since $P$ is 2-Sylow in $G$. Thus $P \cap N$ is 2-Sylow in $N$.

(c) For $g$ of order $2^n$, $l_g$ is a product of $m$ cycles of length $2^n$. $\mathrm{sgn}(l_g) = (-1)^m = -1$ since $m$ is odd.

(d) Induction on $n$. Kernel $K$ of $G \to \{\pm 1\}$ has order $2^{n-1} m$ and contains $g^2$ of order $2^{n-1}$. By induction, $K$ has normal subgroup $M$ of order $m$, which is normal in $G$.
\end{solution}

\subsection*{Team}
\begin{exercise}
Problem 1 (20 points). Find all real orthogonal $2 \times 2$ matrices $k$ with the following property: There is an upper triangular $2 \times 2$ real matrix $b$ with all diagonal entries being positive numbers such that $k b$ is a positive definite symmetric matrix.
\end{exercise}
\begin{solution}
\textbf{Analysis:}
Let $M = kb$ be SPD. Then $b = k^T M$. Since $M$ is SPD, $\det(M) > 0$. Since $b$ is upper triangular with positive diagonal, $\det(b) > 0$. Thus $\det(k) = 1$, so $k$ is a rotation:
$k = \begin{pmatrix} \cos \theta & -\sin \theta \\ \sin \theta & \cos \theta \end{pmatrix}$.
The condition that $b = k^T M$ is upper triangular implies $Y = X \tan \theta$.
Diagonal entries of $b$ are $b_{11} = X/\cos \theta$ and $b_{22} = (Z \cos^2 \theta - X \sin^2 \theta)/\cos \theta$.
For $b_{11} > 0$ and $b_{22} > 0$, we need $\cos \theta > 0$ and $Z > X \tan^2 \theta$.
This is possible for any $\theta \in (-\pi/2, \pi/2)$.
\end{solution}

\begin{exercise}
Problem 2 (20 points). For $x \in \mathbb{Z}$ and $k \geq 0$, define the binomial coefficients $\binom{x}{k} = \frac{x(x-1)\cdots(x-k+1)}{k!}$.
(a) (6 points) Show that $x \in \mathbb{Z} \implies \binom{x}{k} \in \mathbb{Z}$.
(b) (6 points) Show that every function $f : \mathbb{Z}_{\geq 0} \to \mathbb{Z}$ can be expressed as $f(x) = \sum_{k=0}^\infty a_k \binom{x}{k}$, where $a_k \in \mathbb{Z}$ are uniquely determined by $f$.
(c) (8 points) Define $\phi_k(x) = \binom{x + \lfloor k/2 \rfloor}{k}$. Show that every function $f : \mathbb{Z} \to \mathbb{Z}$ can be expressed as $f(x) = \sum_{k=0}^\infty a_k \phi_k(x)$, where $a_k \in \mathbb{Z}$ are uniquely determined by $f$.
\end{exercise}
\begin{solution}
(a) Trivial for $x \geq 0$. For $x = -n < 0$, $\binom{-n}{k} = (-1)^k \binom{n+k-1}{k} \in \mathbb{Z}$.
(b) Newton form: $a_k = \Delta^k f(0) \in \mathbb{Z}$.
(c) Ordered evaluation at $0, -1, 1, -2, 2, \dots$ gives a triangular system with 1s on the diagonal.
\end{solution}

\begin{exercise}
Problem 3 (20 points). Let $K$ be the splitting field of $x^4 - x^2 - 1$ over $\mathbb{Q}$.
(a) (10 points) Show that $\mathrm{Gal}(K/\mathbb{Q}) \cong D_4$.
(b) (10 points) Determine the lattice of subfields of $K$.
\end{exercise}
\begin{solution}
(a) Roots are $\pm \sqrt{\frac{1 \pm \sqrt{5}}{2}}$. $K = \mathbb{Q}(\alpha, i)$ where $\alpha = \sqrt{\frac{1+\sqrt{5}}{2}}$. $[K : \mathbb{Q}] = 8$. Automorphisms form $D_4$.
(b) $D_4$ has 8 subgroups. Fixed fields are subfields of degrees 2 and 4.
\end{solution}

\begin{exercise}
Problem 4 (20 points). Let $G$ be a (not necessarily finite) group and let $F$ be a field of characteristic $\neq 2$. Let $V \neq 0$ be an indecomposable finite-dimensional linear representation of $G$ over $F$. Let $R = \mathrm{End}_F(V)^G$ be the ring of $G$-equivariant endomorphisms of $V$.
(a) (5 points) Prove Fitting's Lemma: every element of $R$ is invertible or nilpotent.
(b) (5 points) Deduce $R/I$ is a division algebra where $I$ is the ideal of non-invertible elements.
(c) (5 points) Show that if there exists a $G$-invariant nondegenerate bilinear form on $V$, then $V$ is orthogonal or symplectic.
(d) (5 points) Assume that $F$ is algebraically closed. Deduce that $V$ cannot be both orthogonal and symplectic.
\end{exercise}
\begin{solution}
(a) $V = \mathrm{ker}(f^N) \oplus \mathrm{im}(f^N)$. Indecomposability implies $f$ is nilpotent or invertible.
(b) $I$ is the set of nilpotents, which is the unique maximal ideal.
(c) Decompose the form into symmetric and alternating parts.
(d) $R/I \cong F$. The space of $G$-invariant forms is 1-dimensional.
\end{solution}

\begin{exercise}
Problem 5 (20 points).
(a) (5 points) For a finite group $G$, show $1 = \sum_{i=1}^h \frac{1}{n_i}$ where $n_i = \#\mathrm{Cent}_G(x_i)$.
(b) (10 points) Show there are finitely many groups with $h$ conjugacy classes.
(c) (5 points) Find all groups with $h=3$.
\end{exercise}
\begin{solution}
(a) Class equation $|G| = \sum \frac{|G|}{n_i}$.
(b) $\sum_{i=1}^h \frac{1}{x_i} = 1$ has finitely many integer solutions.
(c) $1 = 1/|G| + 1/n_2 + 1/n_3$. Solutions correspond to $C_3$ and $S_3$.
\end{solution}
